\documentclass{situationist-preprint}

\title{Two Dogmas, Minus One}
\msid{LNG-04-a}
\msfamily{LNG-04}
\msdiscipline{Philosophy of Language}
\mstype{Principal manuscript}
\mskeywords{analyticity; Quine; Two Dogmas; synonymy; revisability; conceptual change; vagueness; constitutive rules}

\begin{document}
\maketitle

\begin{abstract}
Quine argued in \enquote{Two Dogmas of Empiricism} that the distinction between analytic and synthetic truths cannot be drawn: the notions needed to define analyticity form a closed circle, and no statement is immune from revision in the light of experience. The argument has usually been read as showing that the distinction is empty. We argue that it shows something weaker, that the distinction is vague. We define an analytic truth as one whose acceptance is licensed by the rules that constitute the meanings of its words, and show that the notion so defined has clear positive cases, clear negative cases, and a range of borderline cases, the last generated by the absence of a sharp line between revising a belief and changing a meaning. Quine's circularity argument shows that analyticity has no reductive definition outside its family of notions, which is no evidence of emptiness. His argument from revisability shows that some sentences are borderline, and is answered for the clear cases by the observation that their revision is a change of meaning. There are analytic truths.
\end{abstract}

\section{The question}

\citet[A6--7/B10--11]{Kant1787} distinguished judgements whose predicate is contained in the concept of the subject from those that extend our knowledge beyond it. \citet[\S3]{Frege1884} recast the distinction: a truth is analytic if its proof rests only on general logical laws and definitions. The logical empiricists made it the basis of an epistemology: analytic truths, including logic and mathematics, are true in virtue of meaning and known by knowing meanings; all other knowledge is empirical. \citet{Carnap1947,Carnap1952} gave the most developed version, in which analyticity is relative to a language specified by semantic rules or meaning postulates.

\citet{Quine1951} attacked the distinction in two lines of argument. The first is the circularity argument. Analyticity can be defined in terms of synonymy (an analytic sentence is one that becomes a logical truth by substituting synonyms for synonyms); synonymy in terms of definition, but definitions other than stipulations report prior synonymies; in terms of interchangeability \emph{salva veritate}, but only in a language containing a necessity operator, which presupposes analyticity; or in terms of Carnap's semantic rules, but \enquote{semantic rule} is as obscure as \enquote{analytic}. The family of notions is closed under definition, and none is explained from outside. The second line is the argument from holism. Statements face the tribunal of experience as a corporate body, any statement can be held true come what may if adjustments are made elsewhere, and no statement, not even a law of logic, is immune from revision. So there is no class of statements that are confirmed by any experience whatever, which was the empiricist's second characterization of the analytic.

\citet{GriceStrawson1956} replied that a distinction applied with broad agreement to an open-ended range of new cases is not illusory, whatever the difficulties of defining it, and that the circularity Quine found is the circularity of any family of interdefinable notions. \citet{Putnam1962} argued that Quine was right about most of the statements the empiricists had called analytic, which are framework principles of theories, revisable in scientific revolutions, but that there is a small class of trivial analytic truths involving terms with a single criterion, such as \enquote{all bachelors are unmarried}. \citet{Boghossian1996} distinguished a metaphysical notion, truth in virtue of meaning, which he agreed was defective, from an epistemic notion, on which a sentence is analytic if grasp of its meaning suffices for justified belief in it, which he defended. \citet{Williamson2007} argued that the epistemic notion fails too, since for any sentence a competent speaker may reject it on theoretical grounds while understanding it. \citet{Russell2008} defended a version of the metaphysical notion, and \citet{Chalmers2011} argued that Quine's argument from revisability fails against a notion of apriority that allows for conceptual change.

We argue that there are analytic truths and that Quine's arguments show that the analytic--synthetic distinction is vague, not that it is empty.

\section{Framework}

The account of meaning in LNG-01 provides the notion we need. On that account, the meaning of an expression is constituted by the rules that govern it in a practice, as constrained by truth-conditions.

\begin{definition}[Constitutive rule]\label{d:const}
A rule is \emph{constitutive} of an expression's meaning if it belongs to the use that constitutes that meaning, in the sense of LNG-01, so that a practice in which the rule were replaced would be a practice in which the expression meant something else.
\end{definition}

\begin{definition}[Analytic]\label{d:analytic}
A sentence is \emph{analytic} if its acceptance is licensed by rules constitutive of the meanings of its words, without premises. It is \emph{synthetic} if it is not analytic.
\end{definition}

\begin{definition}[Empty, vague]\label{d:vague}
A distinction is \emph{empty} if one of its sides has no members, or if there is no fact of the matter for any case which side it falls on. It is \emph{vague} if it has clear cases on both sides and borderline cases between them.
\end{definition}

Definition~\ref{d:analytic} is epistemic in Boghossian's sense: it does not say that analytic sentences are made true by meanings, but that their acceptance is licensed by the rules that constitute meanings. \enquote{All vixens are female} is true because all vixens are female; what makes it analytic is that the rule from \enquote{vixen} to \enquote{female fox} is constitutive of \enquote{vixen}. The definition agrees with the account of conceptually grounded a priori knowledge in EPI-12, on which such knowledge is derivation from no premises by rules constitutive of the concepts involved.

\begin{principle}[Graded constitution]\label{p:graded}
Rules differ in how far replacing them would change what an expression means. Some changes clearly change meaning, some clearly do not, and between them there is no sharp line.
\end{principle}

\begin{principle}[Paradigms]\label{p:paradigm}
A distinction that competent speakers apply with agreement to clear cases on both sides, and project with agreement to new cases, is not empty.
\end{principle}

Principle~\ref{p:paradigm} is the argument of \citet{GriceStrawson1956}. It does not say that agreement makes a distinction real; it says that agreement in the projection of a distinction to new cases is evidence that speakers are tracking something, and that the burden is on whoever says they track nothing. Principle~\ref{p:graded} is the claim that will make the distinction vague.

\section{The argument}

\begin{proposition}[Positive cases]\label{p:pos}
There are clear cases of analytic sentences.
\end{proposition}

\begin{proof}
Consider \enquote{if it is raining, then it is raining}, \enquote{every vixen is a fox}, and a sentence introducing a stipulated technical term, such as \enquote{a prime pair is a pair of primes differing by two} in a text that has just introduced \enquote{prime pair}. In each case, acceptance is licensed without premises by rules constitutive of the meanings of the words: the introduction and elimination rules for the conditional, which on LNG-01 fix its meaning; the rule from \enquote{vixen} to \enquote{female fox}, which is what speakers are taught when they learn the word and what they would cite in explaining it; and the stipulation, which is the only rule the term has. A practice in which these rules were replaced would be one in which the words meant something else: a speaker who accepted the conditional's introduction rule but rejected its elimination rule would not be using the conditional, and a speaker who applied \enquote{vixen} to male foxes would not be using \enquote{vixen}.
\end{proof}

\begin{proposition}[Negative cases]\label{p:neg}
There are clear cases of synthetic sentences.
\end{proposition}

\begin{proof}
\enquote{There are foxes in Kent}, \enquote{it is raining in Oxford}. No rule constitutive of the meanings of their words licenses accepting them without premises; acceptance requires observation or testimony.
\end{proof}

\begin{proposition}[Borderline cases]\label{p:border}
There are sentences for which it is indeterminate whether they are analytic.
\end{proposition}

\begin{proof}
Consider the series: \enquote{cats are animals}, \enquote{whales are mammals}, \enquote{gold is a metal}, \enquote{kinetic energy is one half of mass times the square of velocity}, \enquote{atoms cannot be split}. Each was or is accepted by competent speakers in a way that does not rest on specific observations, and each has been revised or could intelligibly be revised. When \enquote{atoms cannot be split} was abandoned, physicists continued to use \enquote{atom} for the same things; the revision was a change of belief. When the classical formula for kinetic energy was replaced in relativistic mechanics, it is unclear whether \enquote{kinetic energy} kept its meaning or acquired a new one, since the rule was central to the term's use but other rules connecting it with work and conservation were retained \citep{Putnam1962}. Whether \enquote{cats are animals} would survive the discovery that the things we call cats are remote-controlled robots, as Putnam supposed, is disputed. By Principle~\ref{p:graded}, there is no sharp line between those revisions that change meaning and those that change belief, so for sentences in the middle of the series it is indeterminate whether the rule that licenses them is constitutive, and so whether they are analytic.
\end{proof}

\begin{theorem}\label{t:main}
The analytic--synthetic distinction is vague and not empty.
\end{theorem}

\begin{proof}
By Propositions~\ref{p:pos} and~\ref{p:neg}, both sides have clear members, so the distinction is not empty by Definition~\ref{d:vague}; Principle~\ref{p:paradigm} confirms this, since speakers agree on these cases and project the distinction to new ones, such as newly stipulated terms. By Proposition~\ref{p:border}, it has borderline cases, so it is vague.
\end{proof}

On the view of vagueness defended in MET-18 and LNG-12, a vague predicate has a sharp boundary that cannot be known. If that view is right, \enquote{analytic} has a sharp boundary somewhere in the series of Proposition~\ref{p:border} that no one can locate. Nothing in Theorem~\ref{t:main} depends on this: the theorem concerns the existence of clear cases and the absence of a knowable sharp line, which every theory of vagueness accepts.

\section{What Quine's arguments show}

\paragraph{Circularity.} The circularity argument shows that \enquote{analytic}, \enquote{synonymous}, \enquote{definition} and \enquote{semantic rule} are interdefinable and that none is defined outside the circle. The premise is true, and Definition~\ref{d:analytic} confirms it: it defines analyticity by constitutive rules, which are part of the same family. But a family of notions none of which is reducible to notions outside it is not thereby empty. \enquote{Cause}, \enquote{law}, \enquote{counterfactual} and \enquote{disposition} form such a family, as do \enquote{knowledge}, \enquote{justification} and \enquote{evidence}. The question is whether the family has application, and Propositions~\ref{p:pos} and~\ref{p:neg} answer it. The circularity argument would show emptiness only given the premise that a notion is legitimate only if it is reducible to behavioral or extensional notions, which is the behaviorism about meaning that Quine later stated explicitly \citep{Quine1960} and that LNG-06 and LNG-05 reject: meaning facts are not fixed by stimulus meaning, and they are not reducible to dispositions.

\paragraph{Revisability.} The argument from holism shows that any sentence may be given up in the light of experience. The premise is true for the sentences of Proposition~\ref{p:pos} as for any others: a community could stop accepting \enquote{every vixen is a fox}. But the analytic sentence is given up by changing the rule that licensed it, and on Definition~\ref{d:const} the word then means something else. Revisability is consistent with analyticity if the revision of an analytic sentence is a change of meaning, and Quine's argument has force only if there is no fact as to whether a given revision is a change of meaning or of belief. For the clear cases there is such a fact. For the middle of the series of Proposition~\ref{p:border} there is not, or none we can know, and this is what Principle~\ref{p:graded} records. So the argument from holism shows that some sentences are borderline, not that none is analytic. This is close to the conclusion of \citet{Chalmers2011}, who argues that conceptual change undercuts the inference from revisability to the absence of apriority.

\paragraph{Confirmation.} Quine's second characterization of the analytic was \enquote{confirmed come what may}. Definition~\ref{d:analytic} does not use it, and Theorem~\ref{t:main} does not need it. Analytic sentences are not confirmed by experience at all, since their acceptance is licensed without premises; that experience can prompt a change in the rules that license them is a different matter, and is the revisability just discussed.

\section{Objections and replies}

\begin{objection}[Understanding without assent]
For any purportedly analytic sentence, there are competent speakers who understand it and reject it: a philosopher who denies that every vixen is a vixen because he holds that sentences with empty terms are false, or a logician who rejects modus ponens for some conditionals \citep{McGee1985}. So acceptance of no sentence is required for understanding, and Definition~\ref{d:analytic} picks out nothing \citep{Williamson2007}.
\end{objection}

\begin{reply}
Definition~\ref{d:analytic} does not say that every competent speaker accepts analytic sentences. It says that their acceptance is licensed by rules constitutive of meaning, and these rules are norms of the community's practice, not dispositions of each speaker (LNG-01). The deviant logician knows the rule he rejects and argues against it; he is a participant in the practice proposing a change in its rules, as a chess player may understand the rules of chess and propose to change one. Williamson's cases show that individual assent is not a criterion of analyticity, which we accept. They do not show that the practice has no constitutive rules. Where the deviant's proposal would, if adopted, change what the word means, his rejection is the first step of a meaning change; whether it would is in some cases indeterminate, which is Proposition~\ref{p:border} again.
\end{reply}

\begin{objection}[Truth in virtue of meaning]
The idea that some sentences are true in virtue of meaning is incoherent. \enquote{Copper is copper} is true because copper is copper, not because of what the words mean; meaning determines what a sentence says, and the world determines whether what it says is so \citep{Boghossian1996}.
\end{objection}

\begin{reply}
We agree. Definition~\ref{d:analytic} is not the metaphysical notion. Analytic sentences are true because things are as they say; what distinguishes them is that the rules constituting the meanings of their words suffice to license accepting them. Whether the metaphysical notion can be rehabilitated, as \citet{Russell2008} argues, does not bear on Theorem~\ref{t:main}.
\end{reply}

\begin{objection}[A vague distinction is no distinction]
The logical empiricists needed a sharp distinction to divide knowledge into two kinds. If the distinction is vague, it cannot do that work, and Quine's conclusion, that there is no principled line, stands.
\end{objection}

\begin{reply}
That the distinction cannot do all the work the empiricists asked of it is part of what we conclude: the frameworks of physical theory, which they classified as analytic, fall largely in the borderline region or on the synthetic side. But a vague distinction does work in its clear cases, as the distinction between day and night does. Logic and some conceptual truths are clearly analytic, and the epistemology of EPI-12, on which they are known a priori by derivation by constitutive rules, applies to them. Quine's conclusion was not that the line is vague but that there is no distinction; that conclusion does not follow.
\end{reply}

\begin{objection}[Indeterminacy]
Quine's later argument for the indeterminacy of translation shows that there is no fact about synonymy, and so none about analyticity.
\end{objection}

\begin{reply}
The argument depends on the claim that translation is constrained only by stimulus meaning. LNG-06 argues that it is constrained further, by compositional structure, eligibility and charity, and that the constraints select a unique translation in the standard case. If they do, there are facts about synonymy, and the argument from indeterminacy does not reach analyticity.
\end{reply}

\section{Conclusion}

Analytic truths are those whose acceptance is licensed by the rules constituting the meanings of their words. There are clear cases, such as the laws of logic, the relations of stipulated terms and the definitional connections of ordinary words, and clear cases of synthetic truths, and between them a range of sentences for which there is no knowable fact as to whether their revision would change meaning or belief. Quine's circularity argument shows that analyticity has no reductive definition, and his argument from holism shows that the borderline region exists. Neither shows that the distinction is empty. One dogma of empiricism, the sharp distinction, should be given up; the distinction itself should not.

\begin{thebibliography}{99}\small
\bibitem[Boghossian(1996)]{Boghossian1996} Boghossian, P.~A. (1996). Analyticity reconsidered. \emph{No{\^u}s}, 30(3), 360--391.
\bibitem[Carnap(1947)]{Carnap1947} Carnap, R. (1947). \emph{Meaning and Necessity}. University of Chicago Press.
\bibitem[Carnap(1952)]{Carnap1952} Carnap, R. (1952). Meaning postulates. \emph{Philosophical Studies}, 3(5), 65--73.
\bibitem[Chalmers(2011)]{Chalmers2011} Chalmers, D.~J. (2011). Revisability and conceptual change in \enquote{Two Dogmas of Empiricism}. \emph{Journal of Philosophy}, 108(8), 387--415.
\bibitem[Frege(1884)]{Frege1884} Frege, G. (1884). \emph{Die Grundlagen der Arithmetik}. Koebner.
\bibitem[Grice and Strawson(1956)]{GriceStrawson1956} Grice, H.~P. and Strawson, P.~F. (1956). In defense of a dogma. \emph{Philosophical Review}, 65(2), 141--158.
\bibitem[Kant(1787)]{Kant1787} Kant, I. (1781/1787). \emph{Kritik der reinen Vernunft}. Hartknoch.
\bibitem[McGee(1985)]{McGee1985} McGee, V. (1985). A counterexample to modus ponens. \emph{Journal of Philosophy}, 82(9), 462--471.
\bibitem[Putnam(1962)]{Putnam1962} Putnam, H. (1962). The analytic and the synthetic. In H.~Feigl and G.~Maxwell (eds.), \emph{Minnesota Studies in the Philosophy of Science}, vol.~3. University of Minnesota Press.
\bibitem[Quine(1951)]{Quine1951} Quine, W.~V.~O. (1951). Two dogmas of empiricism. \emph{Philosophical Review}, 60(1), 20--43.
\bibitem[Quine(1960)]{Quine1960} Quine, W.~V.~O. (1960). \emph{Word and Object}. MIT Press.
\bibitem[Russell(2008)]{Russell2008} Russell, G. (2008). \emph{Truth in Virtue of Meaning: A Defence of the Analytic/Synthetic Distinction}. Oxford University Press.
\bibitem[Williamson(2007)]{Williamson2007} Williamson, T. (2007). \emph{The Philosophy of Philosophy}. Blackwell.
\end{thebibliography}

\end{document}
